\chapter{Améliorations futures}

\section{Restriction des ports RPC dynamiques}

\subsection{Pourquoi restreindre ?}

Par défaut, RPC utilise les ports 49152–65535 (16 383 ports). Écrire des règles pare-feu pour 16 383 ports est ingérable.

Solution : restreindre à une plage fixe de 100 ports (60000–60099).

\subsection{Configuration sur les DCs}

\begin{lstlisting}[style=powershell]
# Sur srv-dc01 et srv-dc02
reg add "HKLM\SOFTWARE\Microsoft\Rpc\Internet" `
    /v Ports /t REG_MULTI_SZ /d "60000-60099" /f

reg add "HKLM\SOFTWARE\Microsoft\Rpc\Internet" `
    /v PortsInternetAvailable /t REG_SZ /d "Y" /f

reg add "HKLM\SOFTWARE\Microsoft\Rpc\Internet" `
    /v UseInternetPorts /t REG_SZ /d "Y" /f

# Redémarrer RPC
Restart-Service RpcSs

# Vérifier
reg query "HKLM\SOFTWARE\Microsoft\Rpc\Internet"
\end{lstlisting}

\subsection{Règle pfSense mise à jour}

Remplacer la règle 49152–65535 par :

\begin{lstlisting}
Action       : Pass
Protocol     : TCP
Source       : LAN net
Destination  : DomainControllers
Dest Port    : 60000 – 60099
Description  : RPC dynamique restreint AD
\end{lstlisting}

\section{Accès Internet sécurisé}

\subsection{Problème actuel}

Actuellement, le proxy Squid est optionnel pour l'IT. Le filtrage URL peut être contourné.

\subsection{Amélioration}

\begin{itemize}
    \item \textbf{NAT obligatoire} : Rediriger HTTP/HTTPS de toutes les interfaces (sauf IT) vers Squid
    \item \textbf{Filtrage DNS} : pfBlockerNG pour bloquer les domaines malveillants au niveau DNS
    \item \textbf{Audit proxy} : Centraliser et analyser les logs Squid pour détection de menaces
\end{itemize}

\section{pfBlockerNG — Blocage DNS}

\subsection{Installation}

\begin{lstlisting}
System > Package Manager > Available Packages
  → pfBlockerNG-devel
  → Install
\end{lstlisting}

\subsection{Configuration}

\begin{lstlisting}
Firewall > pfBlockerNG > DNSBL

Enable DNSBL      : [X]
Active Lists      : [X] ads
                    [X] malware
                    [X] phishing
                    [X] trackers
\end{lstlisting}

pfBlockerNG bloquera les requêtes DNS vers ces domaines.

\section{Suricata IDS/IPS}

\subsection{Détection d'attaques Active Directory}

Suricata est un IDS/IPS capable de détecter les attaques courantes :

\begin{itemize}
    \item Golden Ticket (tickets Kerberos forgés)
    \item DCSync (réplication frauduleuse)
    \item Pass-the-Hash (réutilisation de hash NTLM)
    \item Brute force Kerberos
    \item SMB exploits (EternalBlue, etc.)
\end{itemize}

\subsection{Installation}

\begin{lstlisting}
System > Package Manager > Available Packages
  → Suricata
  → Install

Services > Suricata > Interfaces
  → Add LAN interface
  → Enable
  → ET (Emerging Threats) ruleset
\end{lstlisting}

\section{OpenVPN — Accès télétravail}

\subsection{Cas d'usage}

Permet aux employés de se connecter sécurisement au domaine depuis l'extérieur (télétravail).

\subsection{Configuration rapide}

\begin{lstlisting}
VPN > OpenVPN > Wizards > Local User Access

Protocol             : UDP on IPv4
Port                 : 1194
TLS Authentication   : [X]
Local Network        : 192.168.40.0/24
DNS Server 1         : 192.168.40.10
DNS Default Domain   : corpmsi.local
\end{lstlisting}

\subsection{Déploiement client}

\begin{enumerate}
    \item Installer openvpn-client-export package
    \item Exporter le profil .ovpn
    \item Charger dans OpenVPN Connect (client)
    \item Se connecter
    \item Accéder aux partages DFS et serveurs AD via VPN
\end{enumerate}

\section{Capture de paquets pour debug}

\subsection{Capture trafic Kerberos}

\begin{lstlisting}
Diagnostics > Packet Capture

Interface      : LAN
Host Address   : 192.168.40.10
Protocol       : TCP
Port           : 88 (Kerberos)
Packet Count   : 100
→ Start
\end{lstlisting}

Télécharger le fichier .pcap et analyser dans Wireshark.

\section{Monitoring et SIEM}

\subsection{Tableau de bord}

Métriques à surveiller :

\begin{table}[H]
    \centering
    \begin{tabularx}{\textwidth}{|l|p{4cm}|}
        \hline
        \textbf{Métrique} & \textbf{Outil} \\
        \hline
        Trafic temps réel & pfSense Status → Traffic Graph \\
        \hline
        Règles actives & Firewall → Rules (hit counts) \\
        \hline
        Connexions actives & Diagnostics → States \\
        \hline
        Logs pare-feu & Status → Logs → Firewall \\
        \hline
        Baux DHCP & Status → DHCP Leases \\
        \hline
        Logs Proxy & Diagnostics → Command → tail /var/log/squid/ \\
        \hline
        Logs AD & Event Viewer → Windows Logs \\
        \hline
    \end{tabularx}
    \caption{Métriques et outils de monitoring}
    \label{tab:monitoring}
\end{table}

\subsection{Centralisation SIEM}

Pour une infrastructure production :

\begin{itemize}
    \item \textbf{Elk Stack} (Elasticsearch, Logstash, Kibana)
    \item \textbf{Splunk} (coûteux mais très puissant)
    \item \textbf{Graylog} (open-source, léger)
    
    Envoyer les logs de :
    \begin{itemize}
        \item pfSense (syslog)
        \item Windows Event Log (WEF = Windows Event Forwarding)
        \item Squid/SquidGuard (syslog)
    \end{itemize}
\end{itemize}

\section{Haute disponibilité}

\subsection{DC redondant supplémentaire}

Ajouter srv-dc03 pour redondance 3-way :

\begin{itemize}
    \item Réplication complète
    \item Sites AD pour failover géographique
    \item Load balancing DHCP
\end{itemize}

\subsection{NAS haute disponibilité}

Remplacer srv-fichiers01 par une paire en failover :

\begin{itemize}
    \item Réplication DFS entre deux serveurs
    \item Basculement automatique en cas de défaillance
\end{itemize}

\section{Checklist d'améliorations}

\begin{itemize}
    \item[ ] Restreindre ports RPC à 60000–60099
    \item[ ] Installer et configurer pfBlockerNG
    \item[ ] Installer et configurer Suricata
    \item[ ] Configurer OpenVPN pour télétravail
    \item[ ] Mettre en place SIEM centralisé
    \item[ ] Ajouter DC tertiaire (srv-dc03)
    \item[ ] Mettre en place NAS en failover
    \item[ ] Configurer replication DFS sur 2 serveurs
    \item[ ] Audit DNS avancé (query logging)
    \item[ ] Backup AD automatisé
\end{itemize}
